% =====================================================================
\section{Conclusion and future work}\label{sec:conclusion}
% =====================================================================
\subsection{Conclusion}
This paper presented NITSCI, a compact network for five-level diabetic retinopathy grading from colour fundus
photographs, designed for the common situation in which only a few hundred labelled images are available. The model
places a cross-expert head of 1.40\,M parameters on an EfficientNetV2-S backbone. A global expert relates all retinal
regions through self-attention; a lesion expert receives the same features re-weighted by a lesion-attention map that
is learned from grade labels alone; the two experts exchange information through bidirectional cross-attention; and
an image-adaptive gate decides how much each contributes. A classification head and an ordinal regression head are
combined into one continuous grade score, so that the decision respects the order of the ICDR scale.

The second contribution is the evaluation protocol. The model is pretrained on APTOS-2019, fine-tuned on IDRiD with
stratified five-fold cross-validation, and applied as a twenty-member ensemble of fold models and flips. Every
data-dependent decision, including the choice of decision rule, the four grade cut-points and the referral threshold,
is made on out-of-fold predictions. The held-out test set is used once, and every result is reported with a bootstrap
confidence interval; every ablation is compared with the full model by a paired exact test. This protocol avoids the
optimistic bias that arises when thresholds or models are selected on the test set, a frequent problem in small
medical datasets.

Under this protocol NITSCI achieved a five-class accuracy of \res{acc5}\% (95\% CI \res{acc5ci}) with a quadratic
weighted kappa of \res{qwk5}, and a referable-DR AUC of \res{refauc} with sensitivity \res{refsens}\% and specificity
\res{refspec}\% on the IDRiD test set. \authnote{One or two sentences on the most important ablation findings, for
example the effect of APTOS pretraining and of the cross-expert head relative to the plain baseline, stated with their
confidence intervals or $p$-values.} These results indicate that careful use of a larger in-domain dataset, an
architecture that focuses on lesion evidence without needing lesion annotations, and an ordinal decision rule can
make DR grading on small datasets more reliable. They do not establish clinical readiness, which requires external
and prospective validation.

Table~\ref{tab:summary} links each challenge identified in the Introduction to the component of NITSCI that
addresses it and to the experiment that measures its effect.

\begin{table}[!htbp]
\centering\small
\caption{Summary of the challenges addressed, the corresponding design choices, and where their effect is measured.}
\label{tab:summary}
\begin{tabularx}{\linewidth}{@{}>{\raggedright\arraybackslash}p{3.3cm}>{\raggedright\arraybackslash}X>{\raggedright\arraybackslash}p{3.5cm}@{}}
\toprule
Challenge & Design choice in NITSCI & Evidence\\
\midrule
Small, sparse lesion evidence & Lesion-attention map defining a lesion expert; bidirectional cross-attention with a
  global expert; adaptive gate & Ablations w/o lesion map, w/o cross-attention, w/o gate; plain baseline
  (Table~\ref{tab:ablation})\\
Ordered labels & Classification + ordinal regression heads fused into one score; cut-points fitted on OOF data &
  Ablation w/o ordinal head; decision-rule comparison (Table~\ref{tab:infer}); QWK, within-one-grade accuracy\\
Few labelled images & APTOS-2019 pretraining; EMA weights; five-fold ensemble with TTA & Ablation w/o APTOS;
  single model vs.\ ensemble (Table~\ref{tab:infer})\\
Camera and illumination differences & Ben Graham local-contrast normalisation & Ablation CLAHE
  (Table~\ref{tab:ablation})\\
Class imbalance & Inverse-square-root class weights, label smoothing; balanced metrics & Balanced accuracy, macro F1,
  per-grade results (Table~\ref{tab:perclass})\\
Optimistic evaluation & Test set used once; thresholds on OOF data; bootstrap CIs; McNemar tests &
  Tables~\ref{tab:main}--\ref{tab:ablation}\\
\bottomrule
\end{tabularx}
\end{table}

\subsection{Future work}
Five directions follow directly from the limitations identified in Section~\ref{sec:discussion}:
\begin{enumerate}
\item \textbf{Official partition and benchmark comparison.} Re-running the pipeline on the official IDRiD grading
  partition will allow direct comparison with the challenge results \citep{porwal2020idrid}.
\item \textbf{External validation.} Evaluating the trained ensemble without any re-training on independent datasets
  such as Messidor \citep{decenciere2014messidor} and DDR \citep{li2019ddr}, with subgroup analysis by camera and image
  quality, will measure how well the model generalises.
\item \textbf{Validated lesion attention.} The IDRiD lesion masks make it possible to test whether the learned
  lesion-attention map actually highlights lesions, for example with pointing-game accuracy, and to add weak lesion
  supervision when masks are available \citep{zhou2019collaborative}.
\item \textbf{Stronger initialisation.} Replacing the ImageNet-21k backbone with a retinal foundation model
  \citep{zhou2023foundation} may improve data efficiency further, and the cross-expert head can be attached to any
  backbone that produces a spatial feature map.
\item \textbf{Calibrated, uncertainty-aware referral.} Calibrating the referral probability
  \citep{guo2017calibration} and using the disagreement between ensemble members as an uncertainty signal to
  defer difficult cases to a human grader \citep{araujo2020drgraduate,lakshminarayanan2017simple} would make the system
  safer to deploy in a screening workflow.
\end{enumerate}

\subsection{Recommendations for small-dataset grading studies}
Beyond the specific model, our experience suggests four practices for studies that, like this one, work with a few
hundred labelled images:
\begin{itemize}
\item \textbf{Lock the test set first.} Split it off before any experiment and record its indices; tune nothing on it,
  including thresholds and the choice between decision rules.
\item \textbf{Report uncertainty, not only point estimates.} With fewer than one hundred test images, a bootstrap
  interval is often wider than the differences between methods; reporting it prevents over-claiming.
\item \textbf{Compare under identical conditions.} Baselines and ablations should share splits, pre-processing,
  augmentation and inference; numbers copied from papers that used other splits are not comparable.
\item \textbf{Release the pipeline.} Code that regenerates every number, including the test indices and per-image
  predictions, lets reviewers and other groups verify and extend the results.
\end{itemize}
\sectionend{conclusion}
